\scanpage{103}
et \deleted{\unsure{l'|$\xi$}} on est ramené à prouver l'assertion \deleted{pour} séparément pour $k.1$, et $V$ (égalité des discr.). Or sur $k1$ les deux formes \deleted{bil.} \unsure{qua. ass.|qu. ass.} prennent la $=$ valeur $2$, OK. Sur $V$, \deleted{c'est} elles sont égales ! On gagne !

\deleted{\uline{Proposition}}

\uline{NB} Soit $B$ sous-module de $A$ tel que $B\supset1$, alors
\[
B^{\perp(\varphi)}=B^{\perp(\psi)}.
\]
En effet, \deleted{les deux orthogonaux} \deleted{c'est vrai $B^\perp$ est formé des} on a
\[
\varphi(1,u)=\psi(1,u)=\operatorname{Tr}u
\]
donc les orthogonaux de $1$ sont les $=\text{s}$ (c'est $A^{\unsure{s|\leftrightarrow|\natural}}=\{u\in A\mid\operatorname{Tr}u=0\}$), et $B^\perp\subset A^s$. Or si $\deleted{\unsure{x}}\zeta\in A^s$, on a pour \deleted{\unsure{ql|$\varphi$}} \deleted{tout} $\deleted{\unsure{\xi}}u\in A$
\[
\psi(u,\zeta)=-\varphi(u,\zeta)
\]
donc dire que $\zeta$ est orth. à $B$ pour $\varphi$, ou pour $\psi$, est pareil, cqfd.
% archive sheet number: 102
